\section{二次型与正定性：每个方向上的能量是否为正}
\label{sec:03-Linear-Algebra/06-Symmetric-and-Positive-Definite-Matrices/02}

训练停下来了，梯度范数掉到 \(10^{-8}\)。这说明模型找到了一个极小点吗？

不一定。梯度为零的点还可以是极大点，也可以是鞍点——沿某些方向往上走、沿另一些方向往下走。一阶信息在这三种点上给出的读数完全一样，要区分它们，只能看函数离开这个点之后往哪个方向弯。把驻点挪到原点，二阶 Taylor 展开的主导项是

\[
f(x)\approx f(0)+\tfrac12\,x^{\trans}Hx,
\]

其中 \(H\) 是 Hessian。判断驻点身份，等于判断标量函数 \(x^{\trans}Hx\) 在所有非零方向上的符号。

同一个表达式在别处也在等着。沿方向 \(a\) 观察数据得到的方差是 \(a^{\trans}\Sigma a\)；系统离开平衡位置储存的势能是 \(x^{\trans}Kx\)；带协方差加权的距离平方是 \(u^{\trans}\Sigma^{-1}u\)。三个领域问的是同一个问题：给定对称矩阵 \(A\)，函数

\[
q(x)=x^{\trans}Ax
\]

在 \(x\ne0\) 时是否恒为正？这个函数叫二次型，恒为正这件事叫正定。

\subsection{二次型只看得见对称部分}

先把``\(A\) 对称''这个前提落实。任意实方阵可以拆成对称与反对称两块，\(A=\frac{A+A^{\trans}}2+\frac{A-A^{\trans}}2\)。若 \(K^{\trans}=-K\)，则标量 \(x^{\trans}Kx\) 转置后既等于自身，又等于 \(x^{\trans}K^{\trans}x=-x^{\trans}Kx\)，于是恒为零。反对称那一块对二次型没有任何贡献，

\[
x^{\trans}Ax=x^{\trans}\Bigl(\frac{A+A^{\trans}}2\Bigr)x .
\]

讨论正定性时把矩阵取成对称的，不会丢掉任何信息，反倒是把冗余的自由度去掉了。上一节那套谱工具因此可以直接接上。

\subsection{换到特征坐标，形状只剩符号}

代入 \(A=Q\Lambda Q^{\trans}\)，令 \(y=Q^{\trans}x\)。\(Q\) 正交，所以 \(x\ne0\) 当且仅当 \(y\ne0\)，而且

\[
x^{\trans}Ax=y^{\trans}\Lambda y=\lambda_1y_1^2+\cdots+\lambda_ny_n^2 .
\]

交叉项全部消失，剩下 \(n\) 个独立的平方项，各自配一个特征值当系数。到这一步，\(q\) 的一切定性行为都写在特征值的符号上：全正就没有任何方向能让 \(q\) 变负，出现一负就有方向能让它变负，出现零就有方向让它一动不动。

\begin{center}
\begin{tikzpicture}[x=0.6cm,y=0.6cm,font=\small,>=stealth]
  \begin{scope}[rotate=30]
    \draw[blue!65!black] (0,0) ellipse (2.1 and 1.05);
    \draw[blue!65!black] (0,0) ellipse (1.4 and 0.7);
    \draw[blue!65!black] (0,0) ellipse (0.7 and 0.35);
  \end{scope}
  \fill (0,0) circle (1.6pt);
  \node[below,font=\footnotesize] at (0,-2.9) {\(\lambda_1,\lambda_2>0\)：碗};
  \begin{scope}[shift={(6.2,0)}]
    \draw[black!40,dashed] (-2.3,-2.3) -- (2.3,2.3);
    \draw[black!40,dashed] (-2.3,2.3) -- (2.3,-2.3);
    \draw[red!65!black] plot[domain=0.8:2.3,samples=30] (\x,{sqrt(\x*\x-0.64)});
    \draw[red!65!black] plot[domain=0.8:2.3,samples=30] (\x,{-sqrt(\x*\x-0.64)});
    \draw[red!65!black] plot[domain=0.8:2.3,samples=30] (-\x,{sqrt(\x*\x-0.64)});
    \draw[red!65!black] plot[domain=0.8:2.3,samples=30] (-\x,{-sqrt(\x*\x-0.64)});
    \draw[red!65!black] plot[domain=0.8:2.3,samples=30] ({sqrt(\x*\x-0.64)},\x);
    \draw[red!65!black] plot[domain=0.8:2.3,samples=30] ({-sqrt(\x*\x-0.64)},\x);
    \draw[red!65!black] plot[domain=0.8:2.3,samples=30] ({sqrt(\x*\x-0.64)},-\x);
    \draw[red!65!black] plot[domain=0.8:2.3,samples=30] ({-sqrt(\x*\x-0.64)},-\x);
    \fill (0,0) circle (1.6pt);
    \node[below,font=\footnotesize] at (0,-2.9) {一正一负：鞍};
  \end{scope}
  \begin{scope}[shift={(12.4,0)}]
    \draw[black!55] (-2.0,-2.3) -- (-2.0,2.3);
    \draw[black!55] (2.0,-2.3) -- (2.0,2.3);
    \draw[black!55] (-1.0,-2.3) -- (-1.0,2.3);
    \draw[black!55] (1.0,-2.3) -- (1.0,2.3);
    \draw[very thick,black!75] (0,-2.3) -- (0,2.3);
    \draw[->,black!60] (0,1.2) -- (0,2.1);
    \draw[->,black!60] (0,-1.2) -- (0,-2.1);
    \node[below,font=\footnotesize] at (0,-2.9) {一个为零：槽};
  \end{scope}
\end{tikzpicture}

\emph{二次型的等高线只有三种定性形状。左：两个特征值同为正，等高线是一族同心椭圆，原点是唯一最低点，主轴方向由特征向量给出，特征值越大对应半轴越短。中：一正一负，等高线成双曲线，两条虚线是 \(q=0\) 的方向，原点沿一个方向是最低点、沿另一个方向是最高点。右：一个特征值为零，等高线退化成平行直线，沿粗线方向走多远函数值都不变。}
\end{center}

于是判据可以直接写下：\(A\succ0\)（正定）当且仅当所有 \(\lambda_i>0\)；\(A\succeq0\)（半正定）当且仅当所有 \(\lambda_i\ge0\)。两者的差别集中在中间那张图和右边那张图之间：正定矩阵在每个方向都有严格的向上曲率，半正定允许存在完全平坦的方向，而那些方向恰好张成 \(A\) 的零空间。

顺带留意，只检查坐标轴方向是不够的。矩阵

\[
B=\begin{pmatrix}1&2\\2&1\end{pmatrix}
\]

两个对角元都是正的，\(e_1^{\trans}Be_1=e_2^{\trans}Be_2=1>0\)，可它的特征值是 3 和 \(-1\)，取 \(x=(1,-1)^{\trans}\) 就得到 \(x^{\trans}Bx=-2\)。正定是关于全部方向的断言，对角元只报告了其中 \(n\) 个。

\subsection{几条判据说的是同一句话}

除了特征值，常见的判据还有几条：所有顺序主子式（左上角 \(1\times1,2\times2,\ldots\) 各阶子块的行列式）为正；消元时不换行就能得到全正的主元；存在列满秩的 \(B\) 使 \(A=B^{\trans}B\)；存在对角元为正的下三角 \(L\) 使 \(A=LL^{\trans}\)。教材常把它们排成一张清单，看上去像四条互不相干的检查项，其实它们是同一句话的四种记法。

那句话是：\(q\) 能写成若干个线性表达式的平方和，而且这组表达式不留零方向。特征值版本用正交基做这件事，\(q=\sum\lambda_iy_i^2\) 里的 \(y=Q^{\trans}x\)；\(A=B^{\trans}B\) 版本干脆把平方和写成 \(q=\norm{Bx}^2\)，列满秩正是``\(Bx=0\) 只有零解''，也就是不留零方向；\(A=LL^{\trans}\) 版本把这组线性表达式取成三角的，是下一节的主题；顺序主子式与主元版本则是消元过程留下的记录——消元实际上在做逐步配方，第 \(k\) 个主元就是第 \(k\) 个平方项的系数，而顺序主子式之比恰好等于主元。换基的方式不同，能不能写成没有零方向的平方和这件事不变。

用哪一条取决于手上有什么。有特征值就用符号；只想判断而不想解特征问题，就跑一遍 Cholesky，失败即非正定，代价比特征分解低一个量级；理论推导里想快速证明某个矩阵半正定，\(q=\norm{Bx}^2\) 通常一行就够。

\begin{remark}
半正定的判据不能把``顺序主子式大于零''机械地改成``大于等于零''。取 \(A=\diag(0,-1)\)，两个顺序主子式是 \(0\) 和 \(0\)，都满足非负，但 \(e_2^{\trans}Ae_2=-1\)。半正定要求\emph{所有}主子式非负，不只是顺序的那些；实践中直接看特征值更省事。
\end{remark}

\subsection{\texorpdfstring{为什么 \(A^{\trans}A\) 总是半正定}{为什么 A 的转置乘 A 总是半正定}}

对任意矩阵 \(A\)（不必方，不必对称），

\[
x^{\trans}A^{\trans}Ax=\norm{Ax}^2\ge0 .
\]

这就是上面那句话最短的一个实例。什么时候取到严格不等？当且仅当 \(Ax=0\) 只有零解，即 \(A\) 列满秩。

这条把最小二乘的结构解释清楚了。正规方程 \(A^{\trans}A\hat\beta=A^{\trans}b\) 的系数矩阵永远半正定；设计矩阵列满秩时它正定，目标函数严格凸，最优解唯一。列之间存在共线时它只是半正定，目标沿零空间方向完全平坦，无穷多组参数给出同样的拟合值——图上就是刚才第三张``槽''。回归系数``不稳定''的常见成因就在这里，而不在优化算法。相关讨论见 \hyperref[sec:03-Linear-Algebra/03-Orthogonality/03]{``最小二乘：为无解方程寻找最佳近似''}。

协方差矩阵走的是同一条路。对随机向量 \(X\)，\(\Sigma=\E[(X-\mu)(X-\mu)^{\trans}]\)，把常向量 \(a\) 移进期望即得

\[
a^{\trans}\Sigma a=\E\bigl[(a^{\trans}(X-\mu))^2\bigr]=\Var(a^{\trans}X)\ge0 .
\]

二次型在这里的意思就是``沿方向 \(a\) 看数据时的方差''。它只可能是半正定：若某个非零 \(a\) 让 \(a^{\trans}X\) 几乎处处等于常数，那个方向上的方差为零，\(\Sigma\) 奇异。样本量小于维数时这一定发生——\(n\) 个样本张不出 \(d>n\) 维，样本协方差必然有零特征值。很多算法在这种输入上崩掉，不是实现的问题，是数据本身没提供那些方向的信息。

\subsection{特征值不只定符号，还定收敛速度}

正定性是定性的判断，特征值还携带定量信息。考虑

\[
f(x)=\tfrac12x^{\trans}Ax-b^{\trans}x,\qquad A\succ0,
\]

它的梯度是 \(Ax-b\)，Hessian 恒等于 \(A\)，唯一极小点在 \(x^*=A^{-1}b\)。在特征坐标下这个函数彻底解耦成 \(n\) 条互不相干的抛物线，第 \(i\) 条的曲率是 \(\lambda_i\)。梯度下降用同一个步长 \(\alpha\) 同时处理这 \(n\) 条抛物线：第 \(i\) 个坐标每步乘上 \(1-\alpha\lambda_i\)。要全部不发散就得 \(\alpha<2/\lambda_{\max}\)，可这样一来最平的那个坐标每步只收缩 \(1-\alpha\lambda_{\min}\approx1-2\lambda_{\min}/\lambda_{\max}\)。

\begin{center}
\begin{tikzpicture}[x=0.35cm,y=0.35cm,font=\small,>=stealth]
  \draw[blue!60!black] (0,0) ellipse (12.65 and 2.0);
  \draw[blue!60!black] (0,0) ellipse (10 and 1.581);
  \draw[blue!60!black] (0,0) ellipse (7.071 and 1.118);
  \draw[blue!60!black] (0,0) ellipse (4 and 0.632);
  \fill (0,0) circle (2.4pt);
  \node[below,font=\footnotesize] at (0,-2.4) {\(x^*\)};
  \draw[red!70!black,thick] (11,0.9) -- (10.472,-0.828) -- (9.969,0.762) -- (9.490,-0.701)
     -- (9.034,0.645) -- (8.600,-0.593) -- (8.187,0.546) -- (7.794,-0.502);
  \fill[red!70!black] (11,0.9) circle (2.2pt);
  \node[above,font=\footnotesize] at (11,1.1) {起点};
\end{tikzpicture}

\emph{特征值为 1 和 40 的正定二次型，等高线被拉成一条峡谷。负梯度方向几乎垂直于谷底走向，于是迭代在两侧壁之间来回弹，沿谷底方向每步只挪动一点点。条件数 \(\kappa=\lambda_{\max}/\lambda_{\min}=40\) 正是这两个尺度之比。}
\end{center}

比值 \(\kappa=\lambda_{\max}/\lambda_{\min}\) 叫条件数，它同时控制两件事：椭圆有多扁，以及一阶方法要迭代多少步。\(\kappa=1\) 时等高线是正圆，负梯度直指圆心，一步到位；\(\kappa\) 上到几千，梯度下降的轨迹就变成图里那种锯齿。深度网络训练里的``损失下降极慢但没有停''，很多时候就是这张图的高维版本，而不是学习率设错了数量级。定量的下降估计见 \hyperref[sec:09-Convex-Optimization/06-Optimization-Algorithms/01]{``梯度下降真正难在步长''}。

治法也从这张图读得出来：把坐标按 \(\lambda_i\) 重新缩放，峡谷就变回圆。用 \(A^{-1}\) 做这件事就是 Newton 法，用对角近似做就是 Adam 一类自适应方法，用一个便宜的正定矩阵近似 \(A^{-1}\) 就是预条件。它们的共同前提是 \(A\) 正定——不正定的话，``沿这个方向缩放''里的除法会把下降方向翻个面。

\subsection{回到那个停下来的训练}

对一般的二阶可微函数，Hessian 在某点的特征值给出该点各主方向的曲率。全正意味着所有方向都向上弯，该点是严格局部极小，而且函数在附近局部严格凸——这正是 Newton 步和信赖域方法能用的条件，也是二阶方法要求 \(H\succ0\) 的原因。全负是局部极大。一正一负是鞍点，图里中间那张。全部非负但有零，二阶信息不够判断，得看更高阶项：\(f(x)=x^4\) 与 \(f(x)=-x^4\) 在原点的二阶展开一模一样。

凸性与正定的这层关系在 \hyperref[sec:09-Convex-Optimization/03-Convex-Functions/01]{``弦在图像上方意味着什么''}里会反过来讲一遍：二阶可微函数在凸集上凸，当且仅当 Hessian 处处半正定。严格凸对应处处正定，但反过来不成立，\(f(x)=x^4\) 严格凸而 \(f''(0)=0\)。判据的方向要记准。

高维里还有一个数数的事实值得记住：驻点是极小点，要求 \(d\) 个特征值同时为正。若把这些符号粗略当成独立的，概率是 \(2^{-d}\) 量级。参数上百万时，随手撞见的驻点几乎必然是鞍点，\hyperref[sec:02-Calculus/06-Multivariable-Calculus/05]{``Hessian 与无约束极值''}把这套判据完整建立了一遍。所以开头那个问题的答案通常是：梯度为零不代表到了谷底，得看曲率。

正定性到这里已经完全由特征值刻画了。可是求特征值本身要迭代，要 \(O(n^3)\) 的稠密运算，而且我们经常并不需要知道特征值是多少，只想确认``它们都为正''，或者想把 \(q=\sum\lambda_iy_i^2\) 那个平方和真的写出来用。下一节换一组不正交但三角的基，用一次消元同时完成这两件事。
